% =====================================================================
%  Sections 2.1 - 2.2 : modulations, spectral efficiency, frequencies
% =====================================================================

% ---------------------------------------------------------------------
\subsection{Modulations and spectral efficiency}
\label{sec:2_1}
% ---------------------------------------------------------------------
\paragraph{Data.} $R_b = \SI{30}{\giga\bit\per\second}$, square pulses, WDM grid of \SI{12.5}{\giga\hertz}.

\paragraph{1. Baseband spectral occupancy.}
A modulation format with $M$ symbols carries $n_\mathrm{bit} = \log_2 M$ bits per symbol, so the baud rate
(number of symbols per second) is $D = 1/T_s = R_b/n_\mathrm{bit}$. With rectangular (square) pulses of duration
$T_s$ the PSD is a $\mathrm{sinc}^2$; defining the occupied bandwidth as the first null of the PSD, the baseband
occupancy (on the positive frequency axis) is
\begin{equation}
  B \simeq \frac{1}{T_s} = D = \frac{R_b}{n_\mathrm{bit}} .
\end{equation}
(This is twice the theoretical minimum $B_\mathrm{min} = D/2$ obtained with Nyquist sinc-like pulses.)
\begin{itemize}
  \item \textbf{PAM-2} ($n_\mathrm{bit}=1$): $B = \SI{30}{\giga\bit\per\second}/1 = \SI{30}{\giga\hertz}$.
  \item \textbf{PAM-4} ($n_\mathrm{bit}=2$): $B = 30/2 = \SI{15}{\giga\hertz}$, i.e.\ equal to the baud rate
        $D = R_b/2$.
  \item \textbf{QPSK}: with QAM it is possible to encode two baseband digital transmission signals on the
        same quasi-sinusoidal carrier, using the in-phase (I) and quadrature (Q) components.
        QPSK can be described as two PAM-2 signals in quadrature, each carrying \SI{15}{\giga\bit\per\second},
        so $B = \SI{15}{\giga\hertz}$, the same as PAM-4.
  \item \textbf{16-QAM}: it exploits a PAM-4 on each quadrature. Each branch therefore has the same structure as
        a PAM-4 signal, but carries only half of the bits (\SI{15}{\giga\bit\per\second}), so
        $B = \SI{15}{\giga\hertz}/2 = \SI{7.5}{\giga\hertz}$.
\end{itemize}
\begin{center}
\begin{tabular}{lcccc}
  \toprule
  Format & $n_\mathrm{bit}$ & $D$ [GBd] & $B$ (baseband) & Spectral efficiency $R_b/D$ \\
  \midrule
  PAM-2  & 1 & 30  & \SI{30}{\giga\hertz}  & \SI{1}{bit/s/Hz} \\
  PAM-4  & 2 & 15  & \SI{15}{\giga\hertz}  & \SI{2}{bit/s/Hz} \\
  QPSK   & 2 & 15  & \SI{15}{\giga\hertz}  & \SI{2}{bit/s/Hz} \\
  16-QAM & 4 & 7.5 & \SI{7.5}{\giga\hertz} & \SI{4}{bit/s/Hz} \\
  \bottomrule
\end{tabular}
\end{center}
PAM-4 and QPSK have the same occupancy because they carry the same number of bits per symbol; the advantage
of QPSK is the larger distance between constellation points (better noise tolerance), not the bandwidth.
Increasing the number of bits per symbol reduces the occupancy, at the cost of a higher required SNR.

\paragraph{2. WDM slots.}
Once the real baseband signal of bandwidth $B$ modulates the optical carrier, the resulting bandpass signal
occupies $2B$ around $f_0$: with square pulses the main lobe extends from $f_0 - D$ to $f_0 + D$, so the occupied
optical bandwidth is $B_\mathrm{opt} = 2D$. The number of adjacent \SI{12.5}{\giga\hertz} slots is
$N = \lceil B_\mathrm{opt}/\SI{12.5}{\giga\hertz}\rceil$, and the allocated bandwidth is $N\times\SI{12.5}{\giga\hertz}$:
\begin{center}
\begin{tabular}{lcccc}
  \toprule
  Format & $B_\mathrm{opt} = 2D$ & $B_\mathrm{opt}/\SI{12.5}{\giga\hertz}$ & Slots $N$ & Allocated bandwidth \\
  \midrule
  PAM-2  & \SI{60}{\giga\hertz} & 4.8 & 5 & \SI{62.5}{\giga\hertz} \\
  PAM-4  & \SI{30}{\giga\hertz} & 2.4 & 3 & \SI{37.5}{\giga\hertz} \\
  QPSK   & \SI{30}{\giga\hertz} & 2.4 & 3 & \SI{37.5}{\giga\hertz} \\
  16-QAM & \SI{15}{\giga\hertz} & 1.2 & 2 & \SI{25}{\giga\hertz} \\
  \bottomrule
\end{tabular}
\end{center}
Because the slots must be allocated as whole units, part of the allocated spectrum is unused
(e.g.\ \SI{2.5}{\giga\hertz} for PAM-2, \SI{10}{\giga\hertz} for 16-QAM).


% ---------------------------------------------------------------------
\subsection{Frequencies and wavelengths}
\label{sec:2_2}
% ---------------------------------------------------------------------
\paragraph{Central frequencies.}
Using $f = c/\lambda$ with $c = \SI{3e8}{\metre\per\second}$ and the central wavelength of each band:
\begin{align}
  \text{C-band (1530--1565 nm):}\quad & \lambda_0 = \SI{1547.5}{\nano\metre}
  \;\Rightarrow\; f_0 = \frac{c}{\lambda_0} = \SI{193.86}{\tera\hertz},\\
  \text{L-band (1565--1625 nm):}\quad & \lambda_0 = \SI{1595}{\nano\metre}
  \;\Rightarrow\; f_0 = \frac{c}{\lambda_0} = \SI{188.09}{\tera\hertz}.
\end{align}
(Averaging the frequencies of the band edges instead gives \SI{193.89}{\tera\hertz} and \SI{188.15}{\tera\hertz},
practically the same values.)

\paragraph{Spectral occupancy of the \SI{200}{\giga\bit\per\second} signal.}
As in Section~\ref{sec:2_1}, we assume square pulses, so the optical occupancy is
$B_\mathrm{opt} = 2D = 2R_b/n_\mathrm{bit}$, where the factor 2 accounts for the modulation onto the carrier
(both sides of the main lobe around $f_0$):
\begin{align}
  \text{QPSK } (n_\mathrm{bit}=2):\quad & B_\mathrm{opt} = \frac{2\times\SI{200}{\giga\bit\per\second}}{2} = \SI{200}{\giga\hertz}
  \;\Rightarrow\; N = \frac{200}{12.5} = 16 \text{ slots},\\
  \text{16-QAM } (n_\mathrm{bit}=4):\quad & B_\mathrm{opt} = \frac{2\times\SI{200}{\giga\bit\per\second}}{4} = \SI{100}{\giga\hertz}
  \;\Rightarrow\; N = \frac{100}{12.5} = 8 \text{ slots},
\end{align}
where 16-QAM can be seen as a PAM-4 on each quadrature.
With ideal minimum-bandwidth (Nyquist) pulses, the bandpass occupancy would instead be $B_\mathrm{min} = D$, i.e.\
\SI{100}{\giga\hertz} (8 slots) for QPSK and \SI{50}{\giga\hertz} (4 slots) for 16-QAM; in practice the actual
occupancy lies between these two cases.

\paragraph{Width of the contiguous slots in nanometers.}
Differentiating $\lambda = c/f$ gives $|\Delta\lambda| = \lambda^2\Delta f/c$. Evaluating it at the centre of each band:
\begin{align}
  \Delta\lambda_\mathrm{slot}^{C} &= \frac{(\SI{1547.5}{\nano\metre})^2\times\SI{12.5}{\giga\hertz}}{c} = \SI{0.0998}{\nano\metre},\\
  \Delta\lambda_\mathrm{slot}^{L} &= \frac{(\SI{1595}{\nano\metre})^2\times\SI{12.5}{\giga\hertz}}{c} = \SI{0.1060}{\nano\metre}.
\end{align}
\begin{center}
\begin{tabular}{lccc}
  \toprule
  Signal & Allocated bandwidth & C-band (\SI{1547.5}{\nano\metre}) & L-band (\SI{1595}{\nano\metre}) \\
  \midrule
  QPSK, 16 slots   & \SI{200}{\giga\hertz} & \SI{1.60}{\nano\metre} & \SI{1.70}{\nano\metre} \\
  16-QAM, 8 slots  & \SI{100}{\giga\hertz} & \SI{0.80}{\nano\metre} & \SI{0.85}{\nano\metre} \\
  \bottomrule
\end{tabular}
\end{center}
Since $\Delta\lambda\propto\lambda^2$, the same bandwidth in GHz occupies about $6\%$ more nanometers in the
L-band than in the C-band.
